%%%% ijcai27.tex -- HCAI 2026/27 Lab Project Proposal

\documentclass{article}
\pdfpagewidth=8.5in
\pdfpageheight=11in

\usepackage{ijcai27}

\usepackage[T1]{fontenc}
\usepackage{times}
\usepackage{soul}
\usepackage{url}
\usepackage[hidelinks]{hyperref}
\usepackage[small]{caption}
\usepackage{graphicx}
\usepackage{amsmath}
\usepackage{amsthm}
\usepackage{booktabs}
\usepackage{algorithm}
\usepackage{algorithmic}
\usepackage[switch]{lineno}

% Line numbers are only required for blind-review submissions.
%\linenumbers

\urlstyle{same}

% PDF Info Is REQUIRED (leave untouched).
\ifdefined\pdfinfo
\pdfinfo{
/TemplateVersion (IJCAI.2027.0)
}
\else % XeTeX/LuaTeX (local Tectonic build): same metadata via hyperref
\hypersetup{pdfinfo={TemplateVersion={IJCAI.2027.0}}}
\fi

\title{Who Gets Shortlisted? Auditing and Mitigating Demographic Bias\\
in Multi-Agent LLM Hiring Pipelines for the Portuguese Job Market}

\author{
    Jo\~{a}o Oliveira
    \affiliations
    Human-Centered Artificial Intelligence 2026/27, University of Coimbra
}

\begin{document}

\maketitle

\begin{abstract}
LLM agents are starting to be used to screen and rank job applicants. Prior audits show that single LLM calls score otherwise identical CVs differently depending on the candidate's name. This project proposes an audit of a three-agent hiring pipeline (screener, evaluator, ranker) using counterfactual CVs that differ only in perceived gender and origin, in English and European Portuguese, with origin groups that are relevant to the Portuguese job market. It then evaluates two human-centred mitigations, a blinding agent and rubric-based scoring with quoted evidence. Their justifications are supervised through human-in-the-loop and human-on-the-loop oversight. The project measures both the reduction in bias and the cost in ranking quality. It addresses the HCAI topics of trustworthy AI, multi-agent systems, human oversight and transparency.
\end{abstract}

\section{Problem and Motivation}

Recruitment is turning to agentic AI. Frameworks such as CrewAI and LangGraph already ship example agents that match CVs to positions and recommend candidates. These systems chain several LLM steps, and the recruiter usually sees only the final shortlist. LLMs reproduce social biases from their training data, so two CVs that differ \emph{only} in the candidate's name, gender or nationality can be scored differently. In a pipeline, such a bias may be amplified or masked from one agent to the next without anyone noticing.

The problem matters for three reasons. First, the EU AI Act classifies AI systems for recruitment and candidate selection as \emph{high-risk} (Annex~III), and requires risk management, transparency and effective \emph{human oversight} (Art.~14)~\cite{euaiact2024}. Second, hiring discrimination conflicts with the UNESCO Recommendation on the Ethics of AI~\cite{unesco2021ethics} and with UN SDGs 5, 8 and 10. Third, the risk is real: Amazon abandoned an ML recruiting tool that penalised CVs mentioning women's activities~\cite{dastin2018amazon}.

The problem is also interesting because existing audits mostly test a single LLM call, in English, with US names and racial categories. Little is known about multi-agent pipelines, about non-English use, or about the groups relevant to the Portuguese job market: Portuguese, Brazilian, PALOP (Portuguese-speaking African countries) and Eastern European candidates. The project connects to the course topic of \emph{trustworthy AI}, specifically the requirements of \emph{diversity, non-discrimination and fairness}, \emph{human agency and oversight (HitL/HotL)} and \emph{transparency}~\cite{hleg2019ethics}, and to the topics of \emph{agentic AI} and \emph{AI risks}.

\section{Background}

\paragraph{Hiring discrimination.} Correspondence studies send identical CVs that differ only in the name and measure the gaps in callbacks~\cite{bertrand2004emily}. Blind auditions increased the hiring of women in orchestras~\cite{goldin2000orchestrating}, which motivates \emph{blinding} as a mitigation.

\paragraph{Bias in LLM-based hiring.} Recent audits apply this design to LLMs. They report name-based differences in CV scoring and generation~\cite{armstrong2024silicon}, in hiring-decision emails, with a sensitivity that depends on the prompt~\cite{an2024discriminate}, in retrieval-based CV screening with intersectional effects~\cite{wilson2024resume}, and in GPT rankings of CVs for several jobs~\cite{yin2024bloomberg}. Recent work uses LLM agents to generate counterfactual CV audits with EU AI Act-aligned fairness metrics, but ranks candidates with an embedding model~\cite{yashwant2026counterfactual}, and shows that a two-agent screening procedure changes who advances~\cite{gao2026agent}; neither measures how demographic bias propagates across the agents of a sequential pipeline. Raghavan et al.~\shortcite{raghavan2020mitigating}, writing before LLMs, show that vendors' de-biasing claims are hard to verify.

\paragraph{Fairness and oversight.} Our name-swap design is inspired by \emph{counterfactual fairness}, under which a decision should not change when only the protected attribute changes~\cite{kusner2017counterfactual}. Group fairness is often checked with selection-rate ratios (the US ``four-fifths rule''~\cite{eeoc1978uniform}). Human-centred AI aims for high automation together with high human control~\cite{shneiderman2022hcai} and responsible design~\cite{dignum2019responsible}. Strong LLMs have been used as judges of other models' outputs, but they have biases of their own~\cite{zheng2023judging}.

\section{Goals}

The project aims to measure whether a multi-agent LLM hiring pipeline treats equally qualified candidates differently by gender and origin, and to test whether human-centred mitigations reduce that bias without harming candidate selection. It addresses four research questions:
\begin{itemize}
    \item[\textbf{RQ1}] Do screening, scoring and ranking differ by gender and origin for otherwise identical CVs, and does this depend on the role and on the language (EN vs.\ PT-PT)?
    \item[\textbf{RQ2}] At which agent does the bias appear, and does it grow or shrink along the pipeline?
    \item[\textbf{RQ3}] Do blinding and rubric-with-evidence scoring reduce the bias, and what does each cost in ranking quality?
    \item[\textbf{RQ4}] Are rubric justifications faithful to the CV, and can an LLM judge supervised by a human detect the unfaithful ones?
\end{itemize}

\paragraph{Success criteria.} \emph{Core:} (i)~a reproducible audit of gender~$\times$~origin over three roles in two languages and two LLMs, with effect sizes and corrected significance tests; (ii)~if bias is found, at least one mitigation significantly reduces it (target: the largest score gap halved, with every selection-rate ratio $\geq 0.8$) while losing at most ${\sim}5$ percentage points of pairwise ranking accuracy. A null result also counts: if the 90\% confidence interval of every gap lies within $\pm 2$ score points, the pipeline is reported as equivalent across groups, and the per-agent, language and oversight (RQ4) findings remain contributions. \emph{Stretch:} the combined condition C3, age and career gap / parenthood, and more roles.

\paragraph{Expected contributions.} The project will provide a per-agent bias audit of an agentic hiring pipeline in European Portuguese, with origin groups relevant to the Portuguese job market, which, to our knowledge, have not been studied. It will also compare two oversight-friendly mitigations on both fairness and ranking quality, and release the data and code openly.

\section{Approach}

\paragraph{System under test.} A three-agent pipeline will be built as a sequential CrewAI crew, modelled on public CrewAI recruitment examples. A \emph{Screener} checks the mandatory requirements of the job ad (pass/reject). An \emph{Evaluator} scores each CV that passes from 0 to 100, with a rationale. A \emph{Ranker} orders a pool of eight candidates and returns a top-3 shortlist. The order of the eight CVs in the prompt is shuffled on every run and recorded, so that position effects are not mistaken for demographic bias.

\paragraph{Counterfactual design.} Each \emph{base CV} receives variants that change only demographic cues: name, pronouns, Portuguese grammatical gender (e.g.\ \emph{Engenheiro/Engenheira}) and nationality details. The core design crosses 2~genders with 4~origin groups (Portuguese, Brazilian, PALOP, Eastern European), giving 8 identities. Each identity uses several names, so that no single name drives the results. Ranking pools contain eight CVs of \emph{equal quality} that differ only in identity, so in a fair system every identity is equally likely to reach the top~3; a few \emph{mixed-quality} pools measure the Ranker's quality.

\paragraph{Mitigations.} Table~\ref{tab:conditions} lists the four conditions. Blinding may leave proxies of identity behind (university, languages, Portuguese grammatical gender), and these will be measured. In C2/C3, an LLM judge checks whether each quoted piece of evidence exists in the CV and supports the score given. Under selective \emph{human-in-the-loop} oversight, decisions with a flagged quote or a score near the shortlist cutoff are escalated to a human reviewer; escalations are counted as the cost of oversight, and a sample is reviewed by hand. Under \emph{human-on-the-loop} oversight, the human audits a random sample and monitors the bias metrics.

\begin{table}[t]
    \centering
    \small
    \begin{tabular}{@{}lp{0.70\columnwidth}@{}}
        \toprule
        Condition & Description \\
        \midrule
        C0 Baseline & Pipeline as described above. \\
        C1 Blinding & A pre-processing agent removes names, gender markers, nationality and age cues before screening. \\
        C2 Rubric & Each job criterion is scored separately and must quote the CV text that supports it (structured output). \\
        C3 Both & C1 + C2 (stretch goal). \\
        \bottomrule
    \end{tabular}
    \caption{Experimental conditions.}
    \label{tab:conditions}
\end{table}

\section{Materials}

\paragraph{Data.} The CVs and job ads will be synthetic, generated by an LLM from fixed templates. The core design uses three roles: software engineer (male-stereotyped), nurse (female-stereotyped) and accountant (neutral). Each role has CVs at three quality levels: strong, average and weak. About 10 base CVs per role, each with 8 identity variants in 2 languages, give about 480 CV texts (about 4,000 LLM calls per model). Because the quality level is set when a CV is generated, it serves as \emph{ground truth} for ranking quality. All CVs, including weak ones, meet the mandatory requirements of the job ad, so that every quality level reaches the Evaluator; a few borderline requirements still give the Screener real decisions to make. A sample of the CVs will be checked by hand for realism. Names will be drawn from public name-frequency sources for each group. No real personal data is used. All candidates state the right to work in Portugal and the same Portuguese proficiency, so that origin is not confounded with legal or language requirements, and a manipulation check confirms that the names signal the intended gender and origin. Names only signal \emph{perceived} gender and origin, which is a stated limitation.

\paragraph{Tools.} The project uses Python, CrewAI, and LLMs accessed through OpenRouter: one commercial and one open-weight model, chosen after a pilot. A model from a third family generates the CVs and acts as the judge, so that no model rates its own output. All agents run at temperature~0 with a fixed seed where supported, and a small sample is repeated to check stability. The analysis uses pandas, SciPy and statsmodels, and the figures use matplotlib. The exact models and costs will be reported in the final paper. The code, data and notebooks will be released on GitHub with a master script that reproduces all results.

\section{Evaluation}

\paragraph{Fairness (RQ1--3).} Score gaps are measured as paired differences between variants of the same base CV. They are tested with a mixed-effects model,
$\mathit{score} \sim \mathit{gender} \times \mathit{origin} \times \mathit{cond} + (\mathit{gender} + \mathit{origin}) \times (\mathit{role} + \mathit{lang}) + \mathit{model} + (1 \mid \mathit{CV}) + (1 \mid \mathit{name})$,
with bootstrap confidence intervals and a logistic counterpart for screening and top-3 selection.
Selection rates are computed per identity, both for passing the Screener and for reaching the top~3. Each group's \emph{impact ratio} relative to the most-selected group is compared with the 0.8 threshold. The mean rank per identity in equal-quality pools is tested against a uniform distribution with a permutation test, and the effect of input position is reported separately. All metrics are repeated after each agent to answer RQ2. Multiple comparisons are corrected with the Holm method, and effect sizes are reported.

\paragraph{Ranking quality (RQ3).} Quality is measured as the Spearman correlation between Evaluator scores and the true quality level, and as the Ranker's pairwise accuracy (strong ranked above weak) on the mixed-quality pools, comparing C0 with C1 and C2.

\paragraph{Oversight (RQ4).} The evaluation reports the share of justifications the judge flags as unfaithful, the number of cases escalated to the human (the cost of oversight), and Cohen's $\kappa$ between the judge and two human raters on 50 justifications stratified by the judge's flag, with the human--human $\kappa$ as a reference.

\paragraph{Expected results and plan.} Expected figures include forest plots of score gaps by identity, role and language, impact-ratio heatmaps per agent, and a fairness vs.\ quality trade-off plot per condition. Milestones: CV generation and a pilot of the baseline pipeline by 15~Nov, core experiments by 19~Nov, analysis and human audit by 22~Nov, a paper draft by 24~Nov, and the final paper and code by 27~Nov.

%% The file named.bst is a bibliography style file for BibTeX 0.99c
\bibliographystyle{named}
\bibliography{ijcai27}

\end{document}
